\section{Méthodologie expérimentale}

\subsection{Environnement expérimental}

\begin{table}[H]
    \centering
    \begin{tabular}{l l}
        \toprule
        \textbf{Paramètre} & \textbf{Valeur} \\
        \midrule
        CPU & AMD Ryzen 7 5800H \\
        Cœurs physiques & 8 \\
        Threads logiques & 16 \\
        Python & 3.12.3 \\
        Numba & 0.64.0 \\
        Jeu de données principal & \texttt{data/galaxy\_5000} \\
        Pas de temps & 0.0015 \\
        Grille & \(15 \times 15 \times 1\) \\
        Warmup benchmark & 1 pas \\
        Pas mesurés & 30 \\
        \bottomrule
    \end{tabular}
    \caption{Environnement et paramètres de mesure pour l'étape 1}
    \label{tab:environment_stage1}
\end{table}

\subsection{Procédure}

La procédure de mesure retenue est la suivante :
\begin{enumerate}[label=\arabic*.]
    \item valider le bon fonctionnement de la simulation interactive sur \texttt{data/galaxy\_1000} ;
    \item relever un échantillon des temps \texttt{Render time} et \texttt{Update time} ;
    \item exécuter le benchmark sur \texttt{data/galaxy\_5000} avec \(1, 2, 4, 8, 16\) threads ;
    \item calculer pour chaque cas le speed-up \(S(n) = T(1) / T(n)\) ;
    \item calculer l'efficacité \(E(n) = S(n) / n\).
\end{enumerate}

La campagne de benchmark de cette étape a été exécutée avec \texttt{data/galaxy\_5000}, \(dt = 0.0015\), une grille \(15 \times 15 \times 1\), un pas de warmup non mesuré et 30 pas mesurés par configuration.
